

\section*{Appendix: Extracted Passages from Earlier Drafts}
\textbf{Note from AI:} \textit{The following passages were found in earlier drafts (feb08 and home) but were missing or significantly different in the mar08 version. They have been collected here so you can review and integrate them into the main text as needed.}

\subsection*{Unique Passages from the 'home' Draft}
Network models of cultural transmission and acquisition assume that current tastes are largely a function of the individual’s existing and previous network environment. Cultural capital models on the other hand posit that tastes are driven by durable dispositions acquired throughout childhood and into early to middle adolescence. Given the volatility of social relations uncovered in recent longitudinal research on personal networks, the extant network model therefore implies highly malleable tastes over time. The cultural capital model on the other hand, predicts little over-time volatility in cultural preferences. Unfortunately, we know very little about the longitudinal stability of tastes over time and of how taste changes covary with life course events connected with changing network configurations. Using longitudinal panel data over five years, in this paper I attempt to address this gap through an empirical assessment of the diverging implications of the network and cultural capital models. The results show that consistent with the cultural capital model, tastes are relatively stable over time showing relatively limited amounts of volatility over a five-year period. However, consistent with network theories of taste transmission and decay, the cultural taste change that can be observed is primarily attributable to the life course events most clearly associated with personal network turnover. Finally, and in line with models that posit an alternative avenue of conversion of cultural into social capital, “omnivorous” patterns of cultural taste at a previous time-period are shown to help sustain higher than average rates of social interaction and connections to important social foci into the future, even after holding constant other confounding factors.

Thus, contrary to the network account, according to the cultural capital model tastes should not exhibit over-time volatility, but instead show high inter-temporal correlation. Lizardo (2006) uses this stability assumption to suggest that tastes can in fact be seen as probable (reciprocal) drivers of network connections rather than being seen as exclusively determined by those connections. Using two-stage instrumental variables and matching methods on cross-sectional data, Lizardo (2006) uncovers evidence consistent with this reciprocal effect. This is in agreement with dynamic theories of cultural and network evolution, which rather than postulating a one-way avenue of causation going from networks to cultural tastes propose a reciprocal (over-time) two-way casual interaction between the two (Carley 1991; Mark 1998a). However, even in Lizardo’s (2006) study, the stability hypothesis remains an empirically unverified background assumption. Furthermore to reliance on cross-sectional data, the casual effect of cultural taste on network maintenance or even the more widely accepted effect of network turnover on cultural taste change, has not been decisively established.

It is true that network stability is much more substantial in the “core” personal network, with the bulk of the instability concentrated on the periphery (Bidart and Lavenu 2005; Suitor et al. 1997; Wellman et al. 1997). However, most longitudinal network studies show that the stable “core” network tends to be composed of a small set of alters, primarily kin. The periphery of the network includes most of the persons usually referred to as friends. Studies show that there is considerable even among those ties usually considered “close friends” (Bidart and Lavenu 2005). It is possible to revise network theories of taste transmission to suggest that taste are mostly transmitted through the small, relatively stable core and not through the larger, unstable periphery. However, most contemporary network theories of taste do not make this distinction, and assume that taste are as equally likely to be learned from family and friends (Mark 2003). Interestingly, eliminating the unstable and constantly changing periphery as a relational source of taste transmission would be network theories closer to the presumptions of Bourdieu’s (1984) version of cultural capital theory, in which the family figured as the primary social site in which tastes are learned.

Nevertheless it is clear from recent studies on personal communities that the “core” kin network while mobilized for many purposes, including emotional and financial support, is not used for companionship and sociability (Simmel 1949) in any meaningful sense. Most research and theory on culture consumption communities shows that if tastes are exploited in interaction, this is done precisely with those contacts that are usually named as “friends” and which are seen outside the home and the neighborhood (DiMaggio 1987; Fiske 1987; Frith 1998). However, there is nothing inherent in the logic of the network theory of taste transmission that makes it incompatible with current data large rates of network turnover and the liability of newness of social ties. Nevertheless, it is important to note that the social transmission hypothesis coupled with the empirical observation of the instability of social ties, jointly imply the inter-temporal instability of cultural tastes. The reason for this is that as Mark (2003) notes a taste that is not “reinforced” by social network influence through interaction is likely to be “forgotten.” This is a key assumption of the network based model of taste acquisition. Thus, if we retain the network assumption that individual tastes are largely a function of the tastes displayed by the persons current network alters, and if we accept the fact that these alters are constantly change over the adult life course, then we should expect cultural tastes to themselves exhibit high volatility over time:

While there has been little empirical research on the dynamic stability of tastes through time, there is good reason to suppose that they are more stable than most network theory leads us to believe. For instance, most studies of the role of early family and school experiences on cultural capital have shown strong effects of arts participation, education and after-school training during adolescence on adult tastes even after controlling for subsequent educational attainment (Kracman 1996). Smith (1995) shows that musical tastes are developed early in youth and are fairly stable across the life course (Bennett et al. 1999). From this perspective, while the effect of social network change on cultural taste change is not denied, whatever relational (or “network”) effects on taste exist, must operate in the background of high-levels of cross-temporal taste stability and the reciprocal action of cultural taste on social networks:

In the following study, I use the two-wave panel covering the years 1985 and 1990. While there is data available for these same respondents for the year 1995, I restrict myself to the core set of respondents who participated in both of these waves (N=560).[2] . I do this for two reasons. First, it is well know that cultural participation declines over time, so it is advantageous to measure cultural taste with behavioral data before cultural activity begins to tail off. Second, these two waves used a comparable coding on the taste indicators (a slightly different coding was used in 1995) allowing an operationally valid measure of over-time taste change.

Data limitations. While this data provide one of the few existing quantitative sources of information of cultural taste over-time for a large sample of respondents, there are some limitations of the data at hand that deserve mention: first, there are no direct measures of network connectivity (i.e. size of the personal network). The two measures available, frequency of interaction with friends, and number of associational memberships are important components of being sociable but certainly do not represent the notion of social connectivity in its entirety. Second, there are slight changes in the coding of the cultural activity and friendship interaction variables across waves, with the introduction of a slightly more elaborate scheme in 1990 (including a “monthly” category) and changing to never category to a composite “hardly ever/never.” This may compromise inferences made in investigating taste-loss, by leading to an over-estimation of taste change events for instance. I attempted to deal with this issue by counting a taste loss event as one that involved a shift from one of the extreme categories to the other (which did not change across waves) and by not using items that showed a clear effect of the shift in wording. For instance, it is clear that for the item related to “going to the movies” something like this happened, with 15\% of respondents saying that they “never” go to the movies in 1985, but with 84\% saying that they either “hardly ever or never” go to the movies in 1990. A cross-classification of the responses to the same question across time, reveals that this dramatic shift is mostly due to the fact that most of the respondents who answer “seldom” in 1985 (55\%), where overwhelmingly more likely to answer “hardly ever or never” in 1990. The same artificial taste shift occurs for other cultural taste indicators associated with attending events outside the home (going to a sports event, going to see a dramatic performance in a theater). Cultural practices that do not involve going outside of the home were shown to not be affected by the change in wording, and are therefore the items that I use below in the taste stability and taste loss analyses.

Where Y_(ik)=1 if individual (i) reports that he or she “never” or “seldom” engages in cultural activity (k) in 1990 after reporting that he or she did so “very often” in 1985 and Y_(ik)=0 otherwise and M is the number of sociodemographic factors included in the equations. Y_(ik) is thus our measure that there has been a shift in taste for Y_(k) in the 5 year period. Z in (1) is an i x M design matrix of network turnover indicators. The matrix is constructed from each of the sociodemographic and positional variables (x_(ijm)) whose over-time fluctuation (i.e. going from employed in the first time period to retired in the second) is taken to be likely to lead to cultural taste loss The Z matrix is coded according to the following scheme:

Where everything is as in (5), except that now the observation pertains to the individual-time period. This means that each individual yields two observation, and the random effect (ui) is parameterized as being constant across time but varying across individuals and eit is a classic random disturbance that varies across time and across individuals.

The results shown in the first model of the table are consistent with the culture conversion model: I find that being able to display a wide variety of tastes in the initial time period increases the chances of sustaining higher frequencies of social interaction even after holding constant individual level predictors of interaction frequency in the first time period (γ>0, p<0.05). While not definitive evidence of “causality” these results serve to cement in a more straightforward longitudinal framework those reported by Lizardo (2006) using two-stage methods on cross-sectional data, thus serving to more firmly establish the primary conclusion of the culture conversion model: cultural tastes are used to sustain larger and sparse social networks over time. Thus, culturally active individuals are able to use the conversational resources obtained through the enactment of their tastes (what I refer to as “cultural capital” in this context) to better engage in the conversational rituals necessary for the sustenance of existing friendships and the possible formation of new social ties in interaction. Furthermore, we can see in the fourth column of the table having a wide repertoire of taste is not only useful in increasing social connectivity to other persons over time, but that it is also helps in raising the chances of remaining tied to important social foci of interaction, such as voluntary associations (b₂>0, p<0.05). As Feld (1982) has noted these interaction foci are important since they serve as the primary social mechanism that allows the individual to form social connections which serve to integrate her into the larger social system. Because recruitment (and the over-time stability of membership) into this foci is largely a matter of keeping alive the network ties that bind the individual to other members of the group, it is no surprise that whatever resource helps the individual to sustain social connections to other persons will also appear to sustain his or her connection to these larger social entities.

The results show that the cultural taste instability corollary of traditional network accounts appears not to be operative in this data. Instead, taste for music, newspaper reading, books, and sports show high degrees of stability, a result consistent with the assumptions of a cultural-capital model. However, even though the taste-stability corollary is not supported by these data, the network theoretic assumption that cultural taste loss is more likely to be due to events associated with the disruption of personal networks receives strong support. Across all four cultural forms, taste loss events are strongly associated with changes in interaction frequency with friends, connections to important social foci, and disruptions associated with geographic mobility and changes in work status. Finally, a more stringent test of the process of cultural conversion produces strong support for the theory. Even after taking into account both simultaneity and unobserved heterogeneity bias, we find that individuals who command a wide variety of cultural tastes are more likely to sustain higher levels of social interaction with friends and higher levels of connectivity to important interaction settings (such as voluntary association memberships) through the five year period. This effect it is important to underscore appears not to be due to the reciprocal effect of connectivity on cultural taste breadth.

The results presented here have important implications for how we theorize the statics and dynamics of cultural taste formation, transmission and decay. Rather than the generic idea of the individual who is influenced by the tastes of those persons that he or she is connected to, the cultural capital derived notion of relative taste stability points to two important additional processes usually ignored in traditional network-based accounts:

Choice-homophily of network contacts based on pre-existing cultural dispositions: standard network theories of taste transmission draw on the notion of “homophily” (i.e. the empirical generalization that close friends tend to be similar in sociodemographic characteristics) to explain the fact that cultural tastes and practices tend to cluster in certain areas of social space (Mark 2003). From the standard point of view the homophilous social connection based on socio-demographic similarity pre-exists individual patterns of taste, and this explains why cultural practices “lump” in certain areas of social space. That is, the logic of the theory leads tot the argument that because tastes tend to be associated with social background, and because individuals tend to associate with those who are socio-demographically similar, then we should find that cultural practices occupy delimited niches in social space.

Yet, there is sufficient qualitative evidence (Long 2003; Wali et al. 2002)—and as this paper shows, quantitative evidence (see also Cardon and Granjon 2005; Lizardo 2006)—to suggest that just in the very same that tastes are transmitted through pre-existing network ties, new social ties are formed on the basis of pre-existing cultural dispositions.[8] This is the obverse of the traditional network-theoretic emphasis on the formation of new cultural tastes on the basis of preexisting social ties. In fact, given the relative stability of cultural tastes vis a vis network connections, it is likely that this last process may in fact occur more often that the former process (Lizardo 2006). In this respect (and in contradiction to standard network accounts), the “lumpiness” of cultural tastes in social space may be actually instead be produced through an agentic, active process of contact selection and relationship fine-tuning keyed around cultural compability an observation that would be concordant with Burt’s (2000) observation of the “liability of newness” of recently formed ties, and with recent studies of changes in social networks of young adolescents and young adults (Bidart and Degenne 2005; Bidart and Lavenu 2005; Cardon and Granjon 2005). This would make sociodemographic homophily a by-product of cultural compatibility and not the other way around (this cultural compability may of course be itself due to socio-demographic of the family of origin [Bourdieu 1984]). This is an insight that was always implicit in the theoretical logic of dynamic simulations of network formation based on shared cultural knowledge (Carley 1991), but which has been occluded in recent theory and research due to the one sided emphasis put on the socio-structural determination of shared tastes.

This is important, because the dominant way that “social space” has been conceptualized by network and organizational theorists (i.e. Hannan et al. 2003; McPherson 2004) in relation to culture has been—following Blau (1977)—as a multidimensional space whose very axes are composed of variables (occupational prestige, education) which, according to the culture-conversion model, are themselves implicated in the extent to which tastes may be driven by network constraint or the degree to which that constraint is itself loosened through the ability to deploy those tastes to form social relationships that cut-across social dimensions of stratification and differentiation and which conform to their own—culturally produced—expectations regarding the qualities desirable in a social partner (DiMaggio 1987). In this respect, the dimensions of “social space” are not exogenous to the cultural contents that they contain, but these cultural contents may themselves “warp” these dimensions, such that individuals located at some privileged locations in this space find it easier to traverse a large distance in that space and thus connnect to a wide range of others (Erickson 1996; Lin and Dumin 1986). Individuals located in less privileged areas of the space, on the other hand, are socially and culturally distant from most individuals and cannot connect very well across their locally bounded circle (Bryson 1997).

There is plenty evidence that is consistent with this view. For instance Wellman (1979) and Fischer (1982) show that as socio-economic and educational status increases, individuals are much more likely to have role-segmented, specialized relationships in which a clear differentiation is made between those ties designed for sociability and companionship (usually labeled as “friends”) and social ties that are the result of non-negotiable statuses such as kinship ties. Both studies also show that individuals are much more likely to prefer to obtain companionship from voluntary rather than involuntary ties. These are also the social contacts that the culture conversion model predicts are likely to be formed and sustained through pre-existing forms of cultural compatibility. These are also the social relationships that are likely to be sustained through interaction rituals in which the symbolic goods coming from artworlds and cultural industries figure as the primary subject of conversation (Cardon and Granjon 2005; DiMaggio 1987), and the same ones that are continually undergoing high rates of turnover and being constantly “chosen” anew. In essence, for a relatively privileged cultural and economic elite, as social networks become disembedded from place, kinship and ascription, and come to take the form of increasingly unstable and constantly changing “personal communities” taste may come to become more a driver of network relations, in comparison to individuals who are restricted to a more homogenous and locally constrained set of alters. Consistent with this claim, Lizardo (2006) shows, after holding cultural capital constant, the positive effect of education and socio-economic status on social network size is reduced significantly. This has the curious implication that traditional network theory may more accurately describe the process of cultural taste formation among low cultural capital individuals subject to high levels of network constraint (Burt 2005) while the process of culture conversion may be a more accurate representation of the deployment and operation of cultural taste among culturally advantaged elites.

Morgan, David L. 1989. "Adjusting to widowhood: do social networks really make it easier?" The Gerontologist 29:101-107.

Morgan, David L. and Stephen J. March. 1992. "The Impact of Life Events on Networks of Personal Relationships: A Comparison of Widowhood and Caring for a Spouse with Alzheimer's Disease." Journal of Social and Personal Relationships 9:563-584.

+-------------+--------+--------+-------+-------+--------------+---+----+ |             |        |        | 1990  |       |              |   |    | +=============+========+========+=======+=======+==============+===+====+ |             | Very   | Som    | S     | Mo    | Hardly       |   | T  | |             | often  | etimes | eldom | nthly | Ever/Never   |   | ot | |             |        |        |       |       |              |   | al | |             | (-     | (-     | (-0.  | (0    | (0.6281)     |   |    | |             | 0.5484 | 0.4217 | 0140) | .3560 |              |   |    | |             | )      | )      |       | )     |              |   |    | +-------------+--------+--------+-------+-------+--------------+---+----+ | 1985        |        |        |       |       |              |   |    | +-------------+--------+--------+-------+-------+--------------+---+----+ | Very often  | 43     | 103    | 99    | 31    | 6            |   | 2  | | (-1.11 )    |        |        |       |       |              |   | 82 | +-------------+--------+--------+-------+-------+--------------+---+----+ | Sometimes   | 17     | 57     | 102   | 63    | 16           |   | 2  | | (-0.08)     |        |        |       |       |              |   | 55 | +-------------+--------+--------+-------+-------+--------------+---+----+ | Seldom      | 1      | 1      | 8     | 9     | 5            |   | 24 | | (1.19)      |        |        |       |       |              |   |    | +-------------+--------+--------+-------+-------+--------------+---+----+ |             |        |        |       |       |              |   |    | +-------------+--------+--------+-------+-------+--------------+---+----+ | Total       | 61     | 161    | 209   | 103   | 27           |   | 5  | |             |        |        |       |       |              |   | 61 | +-------------+--------+--------+-------+-------+--------------+---+----+

+------------------+-------------------+---+----------------------------+ |                  | Taste Loss        |   | Taste Loss                 | |                  |                   |   |                            | |                  | (Strong)          |   | (Weak)                     | +==================+===================+===+============================+ |                  | Very Often ->     |   | Very Often ->              | |                  | Never             |   | (Never/Seldom)             | +------------------+-------------------+---+----------------------------+ |                  |                   |   |                            | +------------------+-------------------+---+----------------------------+ | Listening to     | 0.36\%             |   | 1.43\%                      | | Music            |                   |   |                            | +------------------+-------------------+---+----------------------------+ | Reading Books    | 1.07\%             |   | 6.42\%                      | +------------------+-------------------+---+----------------------------+ | Reading Paper    | 0.36\%             |   | 2.50\%                      | +------------------+-------------------+---+----------------------------+ | Following a      | 1.43\%             |   | 4.28\%                      | | Sports Team      |                   |   |                            | +------------------+-------------------+---+----------------------------+

Table 3. Discrete time event-history models showing the effects of change in sociodemographic characteristics associated with network turnover on the probability of cultural taste loss, Project Canada Survey 1985-1990.

-------------------------------------------------------------------------                             Listening   Reading     Reading     Following                             to Music    Books       Paper       Sports   ------------------------- ----------- ----------- ----------- -----------   Δ Social Interaction      1.606*      1.841*      1.821*      1.518*

N                         1121^(a)    1122^(a)    1114^(a)    1118^(a)   -------------------------------------------------------------------------

--------------------------------------------------------------------------                                           N     MMean   Std Dev. Min   Max   --------------------------------------- ----- ------- -------- ----- -----   Culture Consumption Scale (1985)        407   3.50    1.53     0     9

Д Family Arrangements                   407   0.14    ---      0     1   --------------------------------------------------------------------------

[1] ^(*)Data for this project were obtained from the Association of Religion Data Archives at http://www.thearda.com. The author bears sole responsibility for the current interpretation, analyses and tabulations of these data.

[5] The nominal employment status from which the binary indicator for changes in work status is constructed from has six categories (1) full time, (2) part time, (3) unemployed, (4) retired, (5) keeping house, and (6) other. The geographical mobility item is based on the region in which respondent reported residing in each wave. The marital status variable has five categories: (1) married, (2) cohabiting, (3) divorced, (4) widowed and (5) never married.

\subsection*{Unique Passages from the 'feb08' Draft}
Network models of cultural transmission and acquisition assume that current tastes are largely a function of the individual’s existing and previous network environment. Cultural capital models on the other hand posit that tastes are driven by durable dispositions acquired throughout childhood and into early to middle adolescence. Given the volatility of social relations uncovered in recent longitudinal research on personal networks, the extant network model therefore implies highly malleable tastes over time. The cultural capital model on the other hand, predicts little over-time volatility in cultural preferences. Unfortunately, we know very little about the longitudinal stability of tastes over time and of how taste-loss covaries with life course events connected with changing network configurations. Using longitudinal panel data over five years, in this paper I attempt to address this gap through an empirical assessment of the diverging implications of the network and cultural capital models. The results show that consistent with the cultural capital model, tastes are relatively stable over time showing relatively limited amounts of volatility over a five-year period. However, consistent with network theories of taste transmission and decay, the cultural taste loss that can be observed is primarily attributable to the life course events most clearly associated with personal network turnover. Finally, and in line with models that posit an alternative avenue of conversion of cultural into social capital, “omnivorous” patterns of cultural taste at a previous time-period are shown to help sustain higher than average rates of social interaction and connections to important social foci into the future, even after holding constant other confounding factors.

$\ln\left( \frac{p_{ik}}{1 - p_{ik}} \right) = \alpha + {\sum_{m = 1}^{M}\gamma}_{m}Z_{im} + \beta X_{ik} + e_{i}$[] (1)

Table 1. Cross-classification of frequency of interaction with friends outside the neighborhood, Project Canada Survey 1985-1990.

+-------------+--------+--------+------+-------+--------------+---+----+ |             |        |        | 1990 |       |              |   |    | +=============+========+========+======+=======+==============+===+====+ |             | Very   | Som    | Se   | Mo    | Hardly       |   | T  | |             | often  | etimes | ldom | nthly | Ever/Never   |   | ot | |             |        |        |      |       |              |   | al | |             | (      | (      | (0.0 | (-0   | (-0.6281)    |   |    | |             | 0.5484 | 0.4217 | 140) | .3560 |              |   |    | |             | )      | )      |      | )     |              |   |    | +-------------+--------+--------+------+-------+--------------+---+----+ | 1985        |        |        |      |       |              |   |    | +-------------+--------+--------+------+-------+--------------+---+----+ | Very often  | 43     | 103    | 99   | 31    | 6            |   | 2  | | (1.11 )     |        |        |      |       |              |   | 82 | +-------------+--------+--------+------+-------+--------------+---+----+ | Sometimes   | 17     | 57     | 102  | 63    | 16           |   | 2  | | (0.08)      |        |        |      |       |              |   | 55 | +-------------+--------+--------+------+-------+--------------+---+----+ | Seldom      | 1      | 1      | 8    | 9     | 5            |   | 24 | | (-1.19)     |        |        |      |       |              |   |    | +-------------+--------+--------+------+-------+--------------+---+----+ |             |        |        |      |       |              |   |    | +-------------+--------+--------+------+-------+--------------+---+----+ | Total       | 61     | 161    | 209  | 103   | 27           |   | 5  | |             |        |        |      |       |              |   | 61 | +-------------+--------+--------+------+-------+--------------+---+----+

Table 2. Proportion of respondents who experienced a taste loss event, Project Canada Survey 1985-1990.

+------------------+-------------------+---+----------------------------+ |                  | Taste Loss        |   | Taste Loss                 | |                  |                   |   |                            | |                  | (Strong)          |   | (Weak)                     | +==================+===================+===+============================+ |                  | Very Often ->     |   | Very Often ->              | |                  | Never             |   | (Never/Seldom)             | +------------------+-------------------+---+----------------------------+ |                  |                   |   |                            | +------------------+-------------------+---+----------------------------+ | Listening to     | 0.36\%             |   | 1.43\%                      | | Music            |                   |   |                            | +------------------+-------------------+---+----------------------------+ | Reading Books    | 1.07\%             |   | 6.42\%                      | +------------------+-------------------+---+----------------------------+ | Reading Paper    | 0.36\%             |   | 2.50\%                      | +------------------+-------------------+---+----------------------------+ | Following a      | 1.43\%             |   | 4.28\%                      | | Sports Team      |                   |   |                            | +------------------+-------------------+---+----------------------------+

Table 3. Discrete time event-history models showing the effects of change in sociodemographic characteristics associated with network turnover on the probability of cultural taste loss, Project Canada Survey 1985-1990.

-------------------------------------------------------------------------                             Listening   Reading     Reading     Following                             to Music    Books       Paper       Sports   ------------------------- ----------- ----------- ----------- -----------   Δ Social Interaction      1.606*      1.841*      1.821*      1.518*

N                         1121^(a)    1122^(a)    1114^(a)    1118^(a)   -------------------------------------------------------------------------

[1] ^(*)Data for this project were obtained from the Association of Religion Data Archives at http://www.thearda.com. The author bears sole responsibility for the current interpretation, analyses and tabulations of these data.

